% libre-compendium
%
% Copyright (C) 2026 Richard Qin
%
% This document is dual-licensed under either:
% - GNU Free Documentation License, Version 1.3 or any later version
%   published by the Free Software Foundation, with no Invariant Sections,
%   no Front-Cover Texts, and no Back-Cover Texts (see COPYING.DOC); or
% - Creative Commons Attribution-ShareAlike 4.0 International License
%   (see LICENSE-CC-BY-SA-4.0).

\chapter{代数系统}

\section{代数系统的定义与类型}

代数结构研究的对象不是某一个具体数集，而是“集合加上运算”所形成的系统。整数、矩阵、函数、同余类、集合的幂集都可以成为代数系统；关键在于指定哪些对象参与运算，以及运算是否在这些对象中封闭。

\begin{definition}[代数系统]\label{algebra-system}
  设 $X$ 是非空集合，$\Omega$ 是定义在 $X$ 上的若干运算组成的非空集合，则有序对
  \[
  (X,\Omega)
  \]
  称为一个代数系统，或简称代数。
\end{definition}

\begin{definition}[$n$ 元运算]\label{nary-operation}
  设 $n\ge 1$。从 $X^n$ 到 $X$ 的映射
  \[
  \omega:X^n\longrightarrow X
  \]
  称为 $X$ 上的一个 $n$ 元运算；也就是说，每个有序 $n$ 元组
  $(x_1,\ldots,x_n)$ 都唯一地对应一个结果
  $\omega(x_1,\ldots,x_n)\in X$。$0$ 元运算约定为指定 $X$ 中的一个元素。
\end{definition}

若 $\Omega=\{\omega_1,\omega_2,\ldots,\omega_k\}$，其中 $\omega_i$ 是 $n_i$ 元运算，则代数系统的类型记为
\[
(n_1,n_2,\ldots,n_k).
\]
$0$ 元运算表示集合中被指定出来的特殊元素，例如整数代数中的 $0$、$1$。一元运算如取负 $x\mapsto -x$，二元运算如加法、乘法、集合交并。

\begin{pitfall}[运算必须先检查封闭性]\label{closure-first}
  讨论代数系统时，封闭性是入口条件。若 $*$ 被称为 $X$ 上的二元运算，则必须有
  \[
  \forall a,b\in X,\quad a*b\in X.
  \]
  例如正整数集 $\N^+$ 对普通减法不封闭，因此 $(\N^+,-)$ 不能作为以减法为二元运算的代数系统。
\end{pitfall}

更精确地说，判断一个规则是否给出集合 $X$ 上的二元运算，需要同时检查两点：
\[
\forall a,b\in X,\quad a*b \text{ 有唯一结果},
\qquad
\forall a,b\in X,\quad a*b\in X.
\]
前者排除“规则不确定”或“某些输入无法运算”的情形，后者排除结果跑出集合的情形。例如普通除法不是 $\R$ 上的二元运算，因为 $a/0$ 无定义；普通加法不是 $\R^\ast=\R-\{0\}$ 上的二元运算，因为 $1+(-1)=0\notin\R^\ast$。

\begin{example}[title={基本模型}]\label{algebra-examples}
  \begin{enumerate}
    \item $(\Z,+,\cdot,-,0,1)$ 是类型 $(2,2,1,0,0)$ 的代数系统。
    \item 若 $S$ 为集合，则 $(\pow(S),\oplus)$ 是代数系统，其中 $\oplus$ 为对称差。
    \item 时钟代数可写为 $(M,\sigma)$，其中 $M=\{1,2,\ldots,m\}$，$\sigma(k)=k+1$，当 $k=m$ 时令 $\sigma(m)=1$。
    \item 设 $X=\{a,b\}$。$X$ 上所有函数对复合运算构成代数系统，运算表可以刻画其结构。
    \item 同余类集合 $\Z_m=\{[0],[1],\ldots,[m-1]\}$ 在模 $m$ 加法下构成代数系统。
  \end{enumerate}
\end{example}

\begin{definition}[常用运算律]\label{ch07-operation-laws}
  设 $*$ 是 $X$ 上的二元运算。
  \begin{enumerate}
    \item 若 $(a*b)*c=a*(b*c)$ 对任意 $a,b,c\in X$ 成立，则 $*$ 满足\emph{结合律}。
    \item 若 $a*b=b*a$ 对任意 $a,b\in X$ 成立，则 $*$ 满足\emph{交换律}。
    \item 若 $a*a=a$ 对任意 $a\in X$ 成立，则 $*$ 满足\emph{幂等律}。
    \item 设 $\circ$ 也是 $X$ 上的二元运算。若
    \[
    a*(b\circ c)=(a*b)\circ(a*c),\qquad
    (b\circ c)*a=(b*a)\circ(c*a),
    \]
    则称 $*$ 对 $\circ$ 满足左、右分配律；二者同时成立时称满足分配律。
  \end{enumerate}
\end{definition}

\section{运算表、特异元素与一般逆元}

有限集合上的运算常用运算表给出。若表的行和列都按元素 $a_1,\ldots,a_n$ 排列，则第 $i$ 行第 $j$ 列填入 $a_i*a_j$。运算表必须只出现原集合中的元素，否则封闭性失败。

\begin{method}[从运算表读取性质]\label{table-reading}
  \begin{enumerate}
    \item 交换律：运算表关于主对角线对称，即 $a_i*a_j=a_j*a_i$。
    \item 幂等律：主对角线元素满足 $a_i*a_i=a_i$。
    \item 左单位元：某一行与表头完全相同；右单位元：某一列与表头完全相同。二者为同一元素时才是单位元。
    \item 左零元：某一行全为该元素；右零元：某一列全为该元素。二者为同一元素时才是零元。
    \item 逆元：若单位元为 $e$，则在 $a$ 所在行、列中寻找与 $e$ 相交的位置。若存在 $b$ 使 $a*b=b*a=e$，则 $b$ 是 $a$ 的逆元。
  \end{enumerate}
\end{method}

\begin{definition}[左、右单位元与左、右零元]\label{one-zero-sided}
  设 $*$ 是 $X$ 上的二元运算。
  \begin{enumerate}
    \item 若 $e_\ell*x=x$ 对所有 $x\in X$ 成立，则 $e_\ell$ 称为左单位元；若 $x*e_r=x$ 对所有 $x\in X$ 成立，则 $e_r$ 称为右单位元。既是左单位元又是右单位元的元素称为单位元。
    \item 若 $\theta_\ell*x=\theta_\ell$ 对所有 $x\in X$ 成立，则 $\theta_\ell$ 称为左零元；若 $x*\theta_r=\theta_r$ 对所有 $x\in X$ 成立，则 $\theta_r$ 称为右零元。既是左零元又是右零元的元素称为零元。
  \end{enumerate}
\end{definition}

\begin{theorem}[单位元与零元的唯一性]\label{one-zero-unique}
  若 $e_\ell$ 是左单位元，$e_r$ 是右单位元，则 $e_\ell=e_r$，并且所得单位元唯一。若 $\theta_\ell$ 是左零元，$\theta_r$ 是右零元，则 $\theta_\ell=\theta_r$，并且所得零元唯一。若 $X$ 至少含两个元素，则同一运算的单位元与零元不能相等。
\end{theorem}

\begin{proof}
  由 $e_r$ 是右单位元得 $e_\ell*e_r=e_\ell$；由 $e_\ell$ 是左单位元得 $e_\ell*e_r=e_r$，故 $e_\ell=e_r$。若 $e,e'$ 都是单位元，则 $e=e*e'=e'$。零元的证明完全类似：$\theta_\ell*\theta_r=\theta_r$ 且 $\theta_\ell*\theta_r=\theta_\ell$。

  若单位元 $e$ 与零元 $\theta$ 相等，则对任意 $x\in X$，
  \[
  x=x*e=x*\theta=\theta,
  \]
  从而 $X$ 只能有一个元素。这与 $X$ 至少含两个元素矛盾。
\end{proof}

\begin{definition}[一般代数系统中的逆元]\label{general-inverse}
  设 $*$ 有单位元 $e$。若 $y_\ell*x=e$，则称 $y_\ell$ 为 $x$ 的左逆元；若 $x*y_r=e$，则称 $y_r$ 为 $x$ 的右逆元。既是左逆元又是右逆元的元素称为 $x$ 的逆元。
\end{definition}

\begin{theorem}[结合律保证逆元唯一]\label{inverse-unique-associative}
  若 $*$ 满足结合律，且 $x$ 有左逆元 $y_\ell$ 和右逆元 $y_r$，则
  \[
  y_\ell=y_r.
  \]
  因而在结合运算下，若元素有逆元，则逆元唯一。
\end{theorem}

\begin{proof}
  由 $y_\ell*x=e$ 与 $x*y_r=e$ 得
  \[
  y_\ell=y_\ell*e=y_\ell*(x*y_r)=(y_\ell*x)*y_r=e*y_r=y_r.
  \]
\end{proof}

\section{同态、同构与同余}

两个代数系统即使来自不同对象，也可能具有相同的运算规律。同态和同构正是用来描述这种结构对应关系的。

\begin{definition}[同态与同构]\label{hom-iso}
  设 $(X,\Omega_X)$ 与 $(Y,\Omega_Y)$ 是同类型代数系统。映射
  $\varphi:X\to Y$ 若对每个 $n$ 元运算 $\omega\in\Omega_X$ 及其在
  $\Omega_Y$ 中对应的运算 $\bar\omega$，都有
  \[
  \varphi\bigl(\omega(x_1,\ldots,x_n)\bigr)
  =\bar\omega\bigl(\varphi(x_1),\ldots,\varphi(x_n)\bigr),
  \]
  则称 $\varphi$ 为同态映射。对 $0$ 元运算，这一条件表示指定常元被映到对应常元。

  当系统各只写出一个二元运算 $*$、$\circ$ 时，上式就是
  \[
  \varphi(a*b)=\varphi(a)\circ\varphi(b),\qquad a,b\in X.
  \]
  若 $\varphi$ 还是双射，则称为同构映射，并称两个代数系统同构，记为
  \[
  (X,*)\cong(Y,\circ).
  \]
\end{definition}

\lcfigure[caption={代数同态交换图},label={fig:homomorphism-square}]{homomorphism-square}

同态条件表示“运算与映射可交换”：先在原系统中运算再映射，等于先映射到目标系统再运算。同构比同态更强，它要求元素一一对应，因此两个系统在结构上没有本质区别，只是元素名称不同。

\begin{example}[title={奇偶映射保持加法和乘法}]\label{parity-hom}
  定义 $\varphi:\Z\to\{0,1\}$ 为
  \[
  \varphi(n)=
  \begin{cases}
  0,&n\text{ 为偶数},\\
  1,&n\text{ 为奇数}.
  \end{cases}
  \]
  在 $\{0,1\}$ 上定义模 $2$ 加法 $\oplus$ 与模 $2$ 乘法 $\odot$，则
  \[
  \varphi(a+b)=\varphi(a)\oplus\varphi(b),\qquad
  \varphi(ab)=\varphi(a)\odot\varphi(b).
  \]
  因而 $\varphi$ 是从 $(\Z,+,\cdot)$ 到 $(\{0,1\},\oplus,\odot)$ 的同态。两个等式分别说明：和的奇偶性由加数奇偶性决定，积为奇数当且仅当两个因子都为奇数。
\end{example}

\begin{example}[title={行列式同态与自同构}]\label{determinant-auto}
  设 $M_n(\R)$ 为全体 $n$ 阶实矩阵，$\det:M_n(\R)\to\R$ 满足
  \[
  \det(AB)=\det(A)\det(B).
  \]
  因而它是从矩阵乘法代数到实数乘法代数的同态；它通常不是同构，因为不同矩阵可有相同行列式。

  另一方面，在加法代数
  \[
  \Q(\sqrt2)=\{a+b\sqrt2\mid a,b\in\Q\}
  \]
  上，映射 $\sigma(a+b\sqrt2)=a-b\sqrt2$ 满足
  \[
  \sigma(x+y)=\sigma(x)+\sigma(y)
  \]
  且为双射，因此是自同构。
\end{example}

\begin{theorem}[同态像与同态复合]\label{hom-image-composition}
  若 $\varphi:(X,\Omega_X)\to(Y,\Omega_Y)$ 是同态，则 $\img\varphi$ 在
  $\Omega_Y$ 的限制下是 $Y$ 的子代数。若
  \[
  \varphi:(X,\Omega_X)\to(Y,\Omega_Y),\qquad
  \psi:(Y,\Omega_Y)\to(Z,\Omega_Z)
  \]
  都是同态，则 $\psi\circ\varphi$ 也是同态。
\end{theorem}

\begin{proof}
  设 $y_i=\varphi(x_i)\in\img\varphi$。对对应的 $n$ 元运算，
  \[
  \bar\omega(y_1,\ldots,y_n)
  =\bar\omega(\varphi(x_1),\ldots,\varphi(x_n))
  =\varphi(\omega(x_1,\ldots,x_n))\in\img\varphi,
  \]
  故像对每个正元运算封闭；$0$ 元运算由同态定义映入对应常元。又
  \[
  (\psi\circ\varphi)(\omega(x_1,\ldots,x_n))
  =\psi\bigl(\bar\omega(\varphi(x_1),\ldots,\varphi(x_n))\bigr)
  \]
  \[
  =\widetilde\omega\bigl((\psi\circ\varphi)(x_1),\ldots,
  (\psi\circ\varphi)(x_n)\bigr),
  \]
  所以复合保持运算。
\end{proof}

\begin{definition}[同余关系]\label{congruence}
  设 $(X,\Omega)$ 是代数系统，$\sim$ 是 $X$ 上的等价关系。若对每个
  $n$ 元运算 $\omega\in\Omega$，
  \[
  x_1\sim y_1,\ldots,x_n\sim y_n
  \Longrightarrow
  \omega(x_1,\ldots,x_n)\sim\omega(y_1,\ldots,y_n),
  \]
  则称 $\sim$ 是该代数系统上的同余关系。特别地，对二元运算 $*$，条件为
  \[
  a_1\sim a_2,\quad b_1\sim b_2
  \Longrightarrow a_1*b_1\sim a_2*b_2.
  \]
\end{definition}

同余关系比一般等价关系多了“与运算相容”的要求。仅有等价类还不够；若等价类代表元改变后运算结果的等价类也可能改变，则商结构中的运算没有定义。

\begin{theorem}[同态诱导同余]\label{hom-induces-congruence}
  设 $\varphi:(X,*)\to(Y,\circ)$ 是同态。定义
  \[
  a\sim b\Longleftrightarrow \varphi(a)=\varphi(b).
  \]
  则 $\sim$ 是 $(X,*)$ 上的同余关系。
\end{theorem}

\begin{proof}
  等价关系的三条性质由等号的自反性、对称性和传递性直接得到。若 $a_1\sim a_2$ 且 $b_1\sim b_2$，则
  \[
  \varphi(a_1*b_1)=\varphi(a_1)\circ\varphi(b_1)
  =\varphi(a_2)\circ\varphi(b_2)
  =\varphi(a_2*b_2),
  \]
  故 $a_1*b_1\sim a_2*b_2$。
\end{proof}

\section{商代数、自然同态、子代数与直接积}

设 $\sim$ 是 $(X,*)$ 上的同余关系。商集 $X/{\sim}$ 的元素是等价类 $[x]$。在商集上定义
\[
[a]\star[b]=[a*b].
\]
同余条件保证这个定义与代表元选择无关，因此 $(X/{\sim},\star)$ 称为 $(X,*)$ 的商代数。

\begin{theorem}[自然同态与同态分解]\label{quotient-hom}
  映射
  \[
  \pi:X\to X/{\sim},\qquad \pi(x)=[x]
  \]
  是同态，称为自然同态。若同余关系由同态 $\varphi:X\to Y$ 诱导，即
  \[
  a\sim b\Longleftrightarrow \varphi(a)=\varphi(b),
  \]
  则 $X/{\sim}$ 与 $\varphi(X)$ 同构。
\end{theorem}

\begin{proof}
  自然同态满足
  \[
  \pi(a*b)=[a*b]=[a]\star[b]=\pi(a)\star\pi(b).
  \]
  对诱导同余，定义 $\psi:X/{\sim}\to\varphi(X)$ 为
  \[
  \psi([x])=\varphi(x).
  \]
  若 $[x]=[y]$，则 $x\sim y$，从而 $\varphi(x)=\varphi(y)$，故 $\psi$ 良定义。它保持运算、满射且单射，因此为同构。
\end{proof}

\begin{definition}[子代数与直接积]\label{subalgebra-product}
  设 $(X,\Omega)$ 是代数系统，$A\subseteq X$ 且 $A\ne\varnothing$。若每个正元运算将 $A$ 的输入仍映入 $A$，并且每个 $0$ 元运算所指定的元素属于 $A$，则 $(A,\Omega|_A)$ 称为 $(X,\Omega)$ 的子代数。

  若 $(X,*)$ 与 $(Y,\circ)$ 同类型，则其直接积定义在 $X\times Y$ 上：
  \[
  (x_1,y_1)\diamond(x_2,y_2)=(x_1*x_2,\;y_1\circ y_2).
  \]
\end{definition}

直接积按分量运算，因此许多性质会从因子代数传到积代数。例如两个因子运算都可交换、可结合、幂等，则积代数中的对应运算也分别可交换、可结合、幂等；若因子分别有单位元 $e_1,e_2$，则积代数单位元为 $(e_1,e_2)$。消去律也按分量保持：两个因子都满足左（右）消去律时，直接积同样满足左（右）消去律。

\begin{theorem}[同态对特异元素的保持]\label{hom-special-elements}
  设 $\varphi:(X,*)\to(Y,\circ)$ 是满同态。若 $e_X$ 是 $X$ 的单位元，则 $\varphi(e_X)$ 是 $\varphi(X)$ 的单位元；若 $\theta_X$ 是 $X$ 的零元，则 $\varphi(\theta_X)$ 是 $\varphi(X)$ 的零元。若 $x$ 在 $X$ 中可逆，则 $\varphi(x)$ 在 $\varphi(X)$ 中可逆，且
  \[
  \varphi(x^{-1})=\varphi(x)^{-1}.
  \]
\end{theorem}

\begin{proof}
  对任意 $y\in\varphi(X)$，取 $x\in X$ 使 $y=\varphi(x)$，则
  \[
  \varphi(e_X)\circ y=\varphi(e_X)\circ\varphi(x)=\varphi(e_X*x)=\varphi(x)=y,
  \]
  右单位同理。零元证明相同。若 $x*x^{-1}=x^{-1}*x=e_X$，则
  \[
  \varphi(x)\circ\varphi(x^{-1})=\varphi(e_X),
  \qquad
  \varphi(x^{-1})\circ\varphi(x)=\varphi(e_X).
  \]
  故 $\varphi(x^{-1})$ 是 $\varphi(x)$ 的逆元。
\end{proof}

\begin{summary}[本节提要]
  代数系统的核心是集合、运算和封闭性。有限运算可由运算表刻画，单位元、零元、逆元都要相对于指定运算讨论。同态保持运算规律，同构表示结构完全相同。同余关系是能与运算相容的等价关系，它允许在等价类上定义商代数；自然同态把原系统投影到商代数。
\end{summary}
